%% Teste específico para documentos de qualificação
%% Verifica formatação específica para projetos de qualificação

\documentclass[qualificacao-mestrado,onehalfspacing,12pt,hyperref]{UnB-PPGT}

%% Metadados para qualificação de mestrado
\titulo{Análise da Eficiência Energética em Sistemas de Transporte Público Elétrico}
\subtitulo{Proposta de Metodologia para Avaliação de Desempenho}
\autor{Mestrando}{Qualificação Teste}
\orientador{Prof. Dr. Orientador Qualificação}{Universidade de Brasília}
\diamesano{10}{11}{2024}
\coordenador{Prof. Dr. Coordenador PPGT}{Universidade de Brasília}

%% Banca de qualificação (mínimo 3 membros)
\membrobancafunc{Prof. Dr. Orientador Qualificação}{Universidade de Brasília}{Orientador}
\membrobancafunc{Prof. Dr. Examinador 1}{Universidade de Brasília}{Examinador}
\membrobancafunc{Prof. Dr. Examinador 2}{Universidade Federal de Goiás}{Examinador}

%% Palavras-chave para qualificação
\palavraschave[Eficiência energética. Transporte público. Veículos elétricos. Qualificação. Mestrado.]{eficiência energética, transporte público, veículos elétricos, metodologia, avaliação}
\keywords[Energy efficiency. Public transportation. Electric vehicles. Qualification. Master's.]{energy efficiency, public transportation, electric vehicles, methodology, evaluation}

\begin{document}

\pretextual

%% Páginas pré-textuais para qualificação (sem ficha catalográfica)
\capa
\folhaderosto
\folhadeaprovacao

%% Resumo para qualificação
\begin{resumo}
Este projeto de dissertação propõe o desenvolvimento de uma metodologia para análise da eficiência energética em sistemas de transporte público elétrico. O objetivo principal é criar um framework de avaliação que permita comparar diferentes tecnologias de veículos elétricos e otimizar o consumo energético em operações de transporte coletivo.

A motivação para esta pesquisa surge da crescente adoção de veículos elétricos no transporte público brasileiro e da necessidade de ferramentas adequadas para avaliação de seu desempenho energético. A metodologia proposta combinará análise de dados operacionais, modelagem matemática e simulação computacional.

O projeto está estruturado em quatro etapas principais: (i) revisão bibliográfica sobre eficiência energética em transporte público; (ii) desenvolvimento da metodologia de avaliação; (iii) coleta e análise de dados de campo; e (iv) validação através de estudo de caso no Distrito Federal.

Os resultados esperados incluem uma metodologia validada para avaliação de eficiência energética, diretrizes para otimização de operações de transporte elétrico, e contribuições para políticas públicas de mobilidade sustentável.

Esta pesquisa contribuirá para o avanço do conhecimento em transporte sustentável e fornecerá ferramentas práticas para gestores de sistemas de transporte público.
\end{resumo}

%% Abstract para qualificação
\begin{abstract}
This dissertation project proposes the development of a methodology for analyzing energy efficiency in electric public transportation systems. The main objective is to create an evaluation framework that allows comparing different electric vehicle technologies and optimizing energy consumption in public transport operations.

The motivation for this research arises from the growing adoption of electric vehicles in Brazilian public transportation and the need for adequate tools to evaluate their energy performance. The proposed methodology will combine operational data analysis, mathematical modeling, and computational simulation.

The project is structured in four main stages: (i) literature review on energy efficiency in public transportation; (ii) development of evaluation methodology; (iii) field data collection and analysis; and (iv) validation through a case study in the Federal District.

Expected results include a validated methodology for energy efficiency evaluation, guidelines for optimizing electric transport operations, and contributions to sustainable mobility public policies.

This research will contribute to advancing knowledge in sustainable transportation and provide practical tools for public transportation system managers.
\end{abstract}

%% Lista de siglas para qualificação
\begin{siglas}
\item[BRT] \textit{Bus Rapid Transit}
\item[DF] Distrito Federal
\item[GPS] \textit{Global Positioning System}
\item[PPGT] Programa de Pós-Graduação em Transportes
\item[SEMOB] Secretaria de Mobilidade
\item[SOC] \textit{State of Charge}
\item[UnB] Universidade de Brasília
\item[VEB] Veículo Elétrico a Bateria
\end{siglas}

\textual

%% Estrutura típica de projeto de qualificação
\chapter{Introdução}
\label{cap:introducao}

Este projeto de dissertação aborda a análise da eficiência energética em sistemas de transporte público elétrico, tema de crescente relevância no contexto da mobilidade urbana sustentável.

\section{Contextualização}
\label{sec:contextualizacao}

A eletrificação do transporte público representa uma tendência mundial para redução de emissões e melhoria da qualidade do ar urbano.

\section{Justificativa}
\label{sec:justificativa}

A necessidade de metodologias adequadas para avaliação de eficiência energética em veículos elétricos justifica o desenvolvimento desta pesquisa.

\section{Problema de Pesquisa}
\label{sec:problema}

Como desenvolver uma metodologia eficaz para análise da eficiência energética em sistemas de transporte público elétrico?

\section{Objetivos}
\label{sec:objetivos}

\subsection{Objetivo Geral}
Desenvolver uma metodologia para análise da eficiência energética em sistemas de transporte público elétrico.

\subsection{Objetivos Específicos}
\begin{itemize}
\item Realizar revisão bibliográfica sobre eficiência energética em transporte público
\item Identificar variáveis relevantes para avaliação de desempenho energético
\item Desenvolver modelo matemático para análise de eficiência
\item Validar a metodologia através de estudo de caso
\end{itemize}

\chapter{Revisão Bibliográfica}
\label{cap:revisao}

Este capítulo apresenta uma revisão da literatura sobre eficiência energética em sistemas de transporte público elétrico.

\section{Transporte Público Elétrico}
\label{sec:transporte}

O transporte público elétrico tem evoluído significativamente nas últimas décadas, com avanços em tecnologia de baterias e sistemas de propulsão.

\section{Eficiência Energética}
\label{sec:eficiencia}

A eficiência energética em veículos elétricos é influenciada por diversos fatores operacionais e ambientais.

\section{Metodologias de Avaliação}
\label{sec:metodologias}

Diferentes metodologias têm sido propostas na literatura para avaliação de eficiência energética em transporte.

\chapter{Metodologia Proposta}
\label{cap:metodologia}

Este capítulo descreve a metodologia proposta para análise da eficiência energética.

\section{Abordagem Metodológica}
\label{sec:abordagem}

A pesquisa adotará uma abordagem quantitativa baseada em análise de dados operacionais.

\section{Coleta de Dados}
\label{sec:coleta}

Os dados serão coletados através de sistemas embarcados em veículos elétricos em operação.

\section{Modelo Matemático}
\label{sec:modelo}

Será desenvolvido um modelo matemático para relacionar variáveis operacionais com eficiência energética.

\section{Validação}
\label{sec:validacao}

A metodologia será validada através de estudo de caso no sistema de transporte público do Distrito Federal.

\chapter{Cronograma e Recursos}
\label{cap:cronograma}

Este capítulo apresenta o cronograma de execução da pesquisa e os recursos necessários.

\section{Cronograma de Atividades}
\label{sec:cronograma-atividades}

A pesquisa será desenvolvida ao longo de 24 meses, divididos em quatro semestres.

\section{Recursos Necessários}
\label{sec:recursos}

Os recursos incluem equipamentos de medição, software de análise e acesso a dados operacionais.

\chapter{Resultados Esperados}
\label{cap:resultados}

Este capítulo descreve os resultados esperados da pesquisa e suas contribuições.

\section{Contribuições Científicas}
\label{sec:contribuicoes}

A pesquisa contribuirá para o avanço do conhecimento em eficiência energética de transporte público.

\section{Aplicações Práticas}
\label{sec:aplicacoes}

Os resultados terão aplicação direta na gestão de sistemas de transporte público elétrico.

\postextual

%% Bibliografia para qualificação
\begin{thebibliography}{5}
\bibitem{electric2023} SILVA, Ana Maria. \textbf{Eficiência Energética em Ônibus Elétricos: Uma Análise Comparativa}. Revista Brasileira de Transportes, v. 31, n. 2, p. 45-62, 2023.

\bibitem{battery2022} SANTOS, Carlos Eduardo; OLIVEIRA, Maria José. \textbf{Battery Management Systems for Electric Buses: A Review}. Transportation Research Part D, v. 108, Article 103301, 2022.

\bibitem{energy2021} COSTA, Roberto Luiz. \textbf{Energy Consumption Modeling in Electric Public Transportation}. Journal of Public Transportation, v. 23, n. 4, p. 123-140, 2021.

\bibitem{sustainable2023} LIMA, Patricia Fernanda. \textbf{Sustainable Urban Mobility: The Role of Electric Public Transport}. Sustainability, v. 15, n. 8, Article 6789, 2023.

\bibitem{optimization2022} PEREIRA, João Carlos; ALMEIDA, Fernanda Silva. \textbf{Optimization of Electric Bus Operations: A Case Study}. Transportation Science, v. 56, n. 3, p. 678-695, 2022.
\end{thebibliography}

\end{document}